% ---------------------------------------------------------------------------------------------
% DRAFT-NOTE COMMANDS (used in all Sections files; defined once here, the first section imported).
% Remove these two commands and all draft notes before hand-in.
%   \draftnote{...}                 prints its content small and grey under a "DRAFT NOTE" label
%   \figplaceholder{caption}{source} framed empty box for a figure that is to be reused
% \providecommand: nothing happens if a command with the same name already exists.
\providecommand{\draftnote}[1]{\par\begingroup\small\color{black!65}\raggedright
    \expandafter\def\expandafter\UrlBreaks\expandafter{\UrlBreaks\do\-}% let \path{...} break at hyphens
    \noindent\textsf{\textbf{DRAFT NOTE}} \textit{(write the text here, then delete this note)}\par
    #1\par\endgroup}
\providecommand{\figplaceholder}[2]{\par\noindent
    \fbox{\parbox{\dimexpr\linewidth-2\fboxsep-2\fboxrule\relax}{\centering
        \vspace{5mm}\textsf{\textbf{FIGURE PLACEHOLDER}}\\[1mm] #1\\[1mm]
        {\footnotesize Source: #2}\vspace{5mm}}}\par}
% ---------------------------------------------------------------------------------------------

\section*{Introduction}
\addcontentsline{toc}{section}{Introduction}
% NOTE: \section* comes from the Overleaf template. \subsection* + \addcontentsline keep the
% subsections un-numbered (plain \subsection would print as 0.1, 0.2, ...).
% Path shorthands used in all section comments are listed in "Report plan.md".
% Outline shortened 2026-11-23: 5 -> 3 subsections, "Report structure" dropped.

\draftnote{%
\textbf{How to use the notes.} Every subsection has a note with three lists: \textbf{Carry over} (facts
with numbers and source), \textbf{Reuse} (figures as framed boxes; the tables are built in Results) and
\textbf{Fix / gap} (what is wrong, out of date or missing). Page numbers are PDF pages.

\textbf{Source codes} (paths from the group repo root; the last five lines are under
\path{3d-modeling-and-simulation/}):
\begin{itemize}
    \item A1: \path{assignment-1-manufacturing-demonstrator-concept/MMA4002_Manufacturing_Demonstrator_Concept.pdf} (Sept 2026)
    \item A2: \path{assignment-2-product-architectures/Deliverable/MMA4002___Assignment_2_Product_Architectures.pdf} (Oct 2026)
    \item A3: \path{assignment-3-lean/MMA4002___Assignment_3_Lean.pdf} (Oct 2026)
    \item A5 deck / script: \path{assignment-5-industrial-review-board-presentation/MMA4002_IRB_Outdoor_Bench_v2.pptx} and \texttt{\detokenize{IRB presentation script.docx}} (9 Nov 2026); MEDIA: \path{deck/media/} in the same folder
    \item VIS / VISD: \texttt{\detokenize{robot-cell/vision report/jig_vision_report.pdf}} (17 Nov) and \path{jig_vision_report_detailed.pdf} (13 Nov)
    \item PB / PB2: \path{assembly-experiment/report/part_b_report.pdf} and \path{part_b2_report.pdf} (22 Nov); EXP: \path{assembly-experiment/README.md}
    \item QA: \path{qa_notes.md} (dated Q\&A, 15--17 Oct)
    \item CT: \path{robot-cell/cycle_times.csv} (RoboDK era, v3--v8); VER: \path{assembly-experiment/results/versions.csv} (v3--v10, own tools, 20 Nov)
    \item CH: \path{robot-cell/tests/vision_lab/config_history.md}; V11: \path{robot-cell/tests/v11_1plank_RESUME.md} (real robot, 20 Nov)
    \item NS: \texttt{\detokenize{robot-cell/Next steps.md}}; TMB: \texttt{\detokenize{robot-cell/TM bridge plan.md}}; AMR: \texttt{\detokenize{robot-cell/AMR delivery plan.md}}
\end{itemize}}

\subsection*{Background and problem statement}
\addcontentsline{toc}{subsection}{Background and problem statement}
% A4 brief item 1 (problem statement/background, details about the product).
% Customer, user, application and customer need (buyer = municipality/landscape architect, user = the
% public; "classy", traditional look). Why a configurable product family, why automated assembly.
% Sources: A1 report sec. 1.1 (incl. Mouratidis 2026 ref.); index.html (configurator).

\draftnote{%
\textbf{Carry over}
\begin{itemize}
    \item ``The product concept chosen is family of `classy' and non-modernistic outdoor furniture with a focus on benches.'' The end user sits on it; the customer pays and cares more about the message than the price. (A1 p.~2)
    \item ``most people questioned preferred traditional architecture over contemporary (Mouratidis 2026)''; hypothesis: the same holds for outdoor furniture. (A1 p.~2, reference p.~7)
    \item Manufacturing question to borrow: ``How can the demonstrator's product architecture enable simpler, repeatable and scalable manufacturing while preserving required customer functions?'' (A2 brief p.~1)
\end{itemize}
\textbf{Reuse}
\figplaceholder{Four AI-generated concept sketches (label them as AI-generated)}{A1 Fig.~1, p.~2}
\textbf{Fix / gap}
\begin{itemize}
    \item A1 says price matters less, A2 p.~4 says ``affordable'': pick one.
    \item No problem statement for the manufacturing side exists yet.
\end{itemize}}

\subsection*{The demonstrator, objective and scope}
\addcontentsline{toc}{subsection}{The demonstrator, objective and scope}
% Demonstrator - what was built in Manulab: scaled (1:7.5) bench, bent flat-bar frame + planks, backrest
% and decoration variants, laser-cut + 3D-printed parts, cobot assembly cell, AMR delivery.
% One overview figure (render or photo of demonstrator + cell).
% Objective - A4 brief item 1 (objective of the report). Course question (A2 brief p.1): how the product
% architecture enables simpler, repeatable and scalable manufacturing while keeping customer functions.
% Scope - in: demonstrator, robot cell, lean model. Out: real-scale production line.
% Sources: A1 report sec. 1.5; bench-cad/ (BenchInterfaceV2, BenchWOBackrestV2);
% A5 deck/media/ (bench_cad_3d.png, cell_iso.jpg).

\draftnote{%
\textbf{Carry over}
\begin{itemize}
    \item A full-size bench is too large for Manulab, so the demonstrator is scaled 1:7.5, 3D-printed and/or laser-cut wood. (A1 p.~4; A5 deck slide~3)
    \item Cell: Omron TM5-900, Robotiq Hand-E, T-slot plate with base A, base B, plank rack and a two-part guide; the AMR is an Omron LD-250 model in the simulation. (A5 script slide~5; VISD p.~1; NS \S0c)
    \item Message: ``the product is designed for the robot, the cell for the product''. (A5 script slide~3)
\end{itemize}
\textbf{Reuse}
\figplaceholder{CAD render of the bench}{\texttt{\detokenize{bench_cad_3d.png}} (MEDIA)}
\figplaceholder{The robot cell, isometric view}{\texttt{\detokenize{cell_iso.jpg}} (MEDIA)}
\figplaceholder{Photos of the printed benches}{A2 Fig.~18 and 20, pp.~12 and 14}
\textbf{Fix / gap}
\begin{itemize}
    \item The objectives are not written anywhere as a list. Candidates the sources can answer: family from few modules; assembly by one cobot without fasteners; one program for simulation and robot; capacity against takt; jig positions found by vision.
    \item State the scope limits: scaled demonstrator, simulation times, AMR delivery planned but not run.
    \item Say what was built and what was only simulated; A2 p.~15 still says ``Final demonstrator: Not yet finished''.
\end{itemize}}

\subsection*{Team and contributions}
\addcontentsline{toc}{subsection}{Team and contributions}
% Course presentation p.13 ("Document individual contributions"); A2 brief p.4 deliverable 8
% (team contribution log). Short table: member - responsibility - contribution.
% Sources: git log, diary/, A1 report sec. 3 (project plan, Jira).

\draftnote{%
\textbf{Carry over}
\begin{itemize}
    \item Authors on A1--A3: Bjørge Skjelbreid Djuve, John Isak Hammer (title pages).
    \item IRB split: John Isak slides 2--4 and ``learned''; Bjørge slides 5--7 and ``next''. (A5 script)
\end{itemize}
\textbf{Reuse}
\begin{itemize}
    \item Nothing to reuse.
\end{itemize}
\textbf{Fix / gap}
\begin{itemize}
    \item A contribution table (Report plan R3) has to be written by hand: A2 has none, \path{diary/} is empty, and git commit counts under-count CAD and lab work.
\end{itemize}}
